\chapter{Related Work} \label{chapter:related_work}
The usage of automatic stage classification tools is an established field of research,
with efforts dating back over half a century. While earlier methods depended on heuristic
algorithms, modern approaches predominantly use machine learning.

One of the first sources from 1970 \cite{FROST197088} used EEG recordings to 
classify sleep stages in real time using specialized hardware. The system 
first preprocessed the signal by amplifying and filtering it, and then 
operated as a so-called \textit{amplitude-weighted, dominant-frequency meter}.
This module evaluated the EEG signal based on its amplitude across three
predefined levels. Depending on which levels were crossed, two different
artificial signals were generated. These were then fed into a resistor-capacitor
integration circuit with a 15-second memory, determining a final voltage
that was used to classify the sleep stage. To identify REM sleep, the device
used an additional EOG channel to detect rapid eye movements.

Similarly, neonatal sleep stage classification attempts date back decades.
An early method from 1987 \cite{HARPER1987379} used 7 cardiac and respiratory measures
of children up to age of 6 months. Using discriminant analysis they developed a model,
for which article reports an overall accuracy of 84.8\%.

However, since the rise of machine learning, more advanced methods have been developed.
Most often, these methods either use EEG or all polysomnography (PSG) signals as input. It is however, most common, 
for these models to learn on a specific demographic. That is, models either see only adult or only
neonatal patients.

\section{Machine Learning Methods for Sleep Stage Classification of Adults}

A random forest classifier was one of the first machine learning models used for sleep stage classification \cite{FRAIWAN201210}.
The model used a single EEG channel represented in the time-frequency domain. Their study consisted of 16 
subjects aged 30 to 75. Researchers tried multiple feature extraction methods and found the 
best results using Continuous Wavelet Transform signal breakdown. Reportedly, 
the model achieved an accuracy of 83\%.

Later, specialized models designed for temporal data were used. A combination of convolutional neural network (CNN)
and long short-term memory (LSTM) network was developed in 2021 \cite{cnnlstm2021}.
Model used three convolutional layers and two LSTM layers and finally a dense layer 
for classification into one of the 5 stages. The model used two EEG channels as input.
The model was tested on 12 subjects, in ages between 21 and 61. The model achieved
an accuracy of 74\%.

Additionally, a transformer-based model called MultiChannelSleepNet was developed \cite{transformersleep2023}.
Rather than operating on raw signals, the model uses time-frequency representations of multiple EEG and EOG
channels: transformer encoders first extract features from each channel independently, and a second, smaller
set of transformer encoders then fuses these into joint features. The model was evaluated on three datasets,
including SleepEDF-78, which contains subjects aged between 25 and 101 years, on which it achieved an accuracy
of 85.0\%.

\section{Machine Learning Methods for Sleep Stage Classification of Infants and Children} 
Similarly to sleep staging for adults, temporal models for neonatal
classification have been developed. Approach from article \cite{neonetalLstm} uses 
a single EEG channel to classify sleep into 3 stages, Awake, Active sleep and Quiet sleep.
All patients were 40 weeks post-conceptional old. Using a bidirectional LSTM network,
and a fully connected classification layer the model achieved an overall accuracy of
96.81\%. Model was evaluated using a 10-fold cross validation \cite{neonetalLstm}.

A CNN model from 2023 aimed to classify sleep 
into Quiet and Non-Quiet stages \cite{neonatalcnn2023}.
The model was trained on EEG data of 19 preterm and term newborns,
using 9 channels as input. The model achieved an accuracy of 94.07\% \cite{neonatalcnn2023}.
The EEG data was first preprocessed into 8 time-domain features and 
four frequency-domain features, which were then used as input into the CNN model.

Similarly to adults, a transformer based model was developed for children sleep staging \cite{transformerinfant2022}.
The architecture used was inspired by Vision Transformers (ViT), which is commonly 
used in image based tasks. However, the architecture was 
modified for multichannel input. An epoch, which consists of 30 seconds of data,
is split into 30 patches, one for each second. This is analogous to tokenization.
From this, the model classifies each epoch into one of the standard 5 stages.
The data consisted of 3.6 million epochs from infants, children and adolescents patient.
Model achieved an overall accuracy of 78.2\%, with 
highest accuracy for patients in the range of 8 to 9 years (85\%) \cite{transformerinfant2022}.